\section{Conclusion}
\label{sec:conclusion}

This Systematic Literature Review synthesized 18 peer-reviewed foundational and empirical studies investigating the effect of prototyping on identifying missing and unclear requirements in software engineering. Structured around the IREB framework and tripartite analytical dimensions, our findings address the three core research questions:

First, regarding \textit{prototyping modalities and process (RQ6.1)}, prototypes operationalized under the PAM framework function as boundary objects across the RE lifecycle: exploratory probes surface tacit knowledge during elicitation, while interactive simulations catalyze requirement evolution and uncover up to 42\% additional missing requirements during validation.

Second, regarding \textit{efficacy compared to textual specifications alone (RQ6.2)}, prototyping significantly outperforms textual specifications and unstructured interviews in gross functional requirement discovery volume ($p = 0.02$) and provides a gender-egalitarian medium ($p = 0.132$). However, passive prototypes do not guarantee specification completeness ($p = 0.17$) or correctness ($p = 0.59$) for non-visual backend logic.

Third, regarding \textit{susceptibility by requirement typology (RQ6.3)}, prototypes excel in clarifying functional workflows and usability/GUI requirements ($p = 0.000$). Traditional non-functional requirement deficits are overcome via augmented modalities, including think-aloud paper prototyping ($p = 0.001$) and contextual video scenarios.

In conclusion, prototyping is an essential, evidence-backed strategy for reducing omissions and linguistic ambiguities. To maximize efficacy while mitigating risks, teams should deploy low-fidelity sketches during exploratory framing, adopt Enriched User Stories with visual acceptance criteria during backlog grooming, and transition to high-fidelity mockups only during confirmatory pre-release validation.
